\documentclass[a4paper,11pt]{article}
\usepackage[utf8]{inputenc}
\usepackage[T1]{fontenc}
\usepackage[lithuanian]{babel}
\usepackage{enumitem}
\usepackage{geometry}
\usepackage{array}
\geometry{margin=2.5cm}

\title{Projekto valdymo planas\\\large Pokyčių valdymo procedūra (Change Management)}

\begin{document}
\maketitle

\section{Tikslas ir taikymo sritis}
Šio dokumento tikslas – apibrėžti formalų ir kontroliuojamą procesą, skirtą bet kokiems projekto apimties, tvarkaraščio, biudžeto ar kokybės reikalavimų pakeitimams valdyti po to, kai pradiniai projekto planai (bazės) yra oficialiai patvirtinti. 

Pokyčių valdymu siekiama:
\begin{itemize}
    \item Apsaugoti projektą nuo nekontroliuojamo apimties išsiplėtimo (\textit{angl. scope creep}).
    \item Užtikrinti skaidrų ir pagrįstą sprendimų priėmimą.
    \item Garantuoti, kad visi suinteresuotieji asmenys laiku sužino apie planuojamų pakeitimų pasekmes.
\end{itemize}

\section{Pokyčių klasifikacija}
Visi projekto pakeitimai skirstomi į tris pagrindines kategorijas pagal jų poveikį projekto trikampiams apribojimams:
\begin{enumerate}
    \item \textbf{Nedideli pakeitimai:} Įtakos neturi nei projekto biudžetui, nei galutiniam terminui (pvz., nedidelis užduočių eiliškumo pasikeitimas komandos viduje).
    \item \textbf{Vidutiniai pakeitimai:} Keičia tarpinius terminus arba reikalauja nedidelio resursų perskirstymo neviršijant bendro biudžeto rezervo.
    \item \textbf{Dideli pakeitimai:} Iš esmės keičia projekto apimtį, reikalauja papildomo finansavimo arba lemia galutinio projekto termino pratęsimą.
\end{enumerate}

\section{Pokyčių valdymo procesas}
Bet koks pokytis projekto eigoje privalo praeiti šiuos 5 nuoseklius etapus:

\begin{enumerate}
    \item \textbf{1 etapas: Pokyčio inicijavimas ir pateikimas}
    Bet kuris projekto komandos narys ar užsakovas gal pastebėti pokyčio poreikį. Iniciatorius privalo užpildyti formalią \textit{Pokyčio užklausą} (angl. \textit{Change Request}), kurioje aiškiai aprašo pokyčio esmę, priežastį ir naudą.
    
    \item \textbf{2 etapas: Registracija}
    Projekto vadovas gauna užklausą, ją užregistruoja centrinėje sistemoje – \textit{Pokyčių registre} (angl. \textit{Change Log}) – ir suteikia unikalų numerį tolimesniam sekančiam sekimui.
    
    \item \textbf{3 etapas: Poveikio analizė}
    Projekto vadovas kartu su atsakingais specialistais įvertina pokyčio poveikį:
    \begin{itemize}
        \item \textit{Apimčiai:} Kokios naujos funkcijos, reikalavimai atsiranda arba yra pašalinami?
        \item \textit{Laikui:} Kiek dienų ar savaičių užtruks įgyvendinimas?
        \item \textit{Kaštams:} Kokia finansinė naštos ar sutaupymo įtaka bus sugeneruota?
        \item \textit{Rizikai:} Kokias naujas grėsmes ar galimybes tai sukuria projekto sėkmei?
    \end{itemize}
    
    \item \textbf{4 etapas: Sprendimo priėmimas}
    Remiantis poveikio analize, pokyčio užklausa teikiama tvirtinti atitinkamam valdymo organui (pagal kompetencijos ribas). Galimi sprendimai: patvirtinti, atmesti arba grąžinti patikslinimui.
    
    \item \textbf{5 etapas: Įgyvendinimas ir atnaujinimas}
    Jei pokytis patvirtinamas, projekto vadovas oficialiai atnaujina projekto planą, biudžetą bei tvarkaraštį ir paveda komandai įgyvendinti užduotį. Atmesti prašymai archyvuojami nurodant atmetimo priežastį.
\end{enumerate}

\section{Kompetencijos ir sprendimų priėmimo ribos}
Lentelėje pateikiama atsakomybės matrica už pokyčių tvirtinimą pagal jų finansinį ir laiko poveikį:

\begin{center}
\begin{tabular}{|l|p{6cm}|l|}
\hline
\textbf{Pokyčio lygis} & \textbf{Kriterijai / Apibūdinimas} & \textbf{Tvirtinantis asmuo} \\ \hline
Žemas & Biudžeto pokytis iki 1,000 € / laiko vėlavimas iki 3 d. & Projekto vadovas \\ \hline
Vidutinis & Biudžeto pokytis 1,000 – 5,000 € / vėlavimas 3–10 d. & Projekto rėmėjas \\ \hline
Aukštas & Biudžeto pokytis virš 5,000 € / esminis apimties keitimas & Pokyčių valdymo komitetas (CCB) \\ \hline
\end{tabular}
\end{center}

\end{document}